% Prose rendering of PrimeTensor/Scale/Laws.lean.
% Fragment: no \documentclass, no preamble, no document environment.

\section{Composition laws for the intrinsic scale}
\label{Scale:Laws:sec}

\leanfile{PrimeTensor/Scale/Laws.lean}

\subsection*{What this file establishes}

\texttt{PrimeTensor/Scale.lean} built the ladder of closeness relations
$\name{ScaleWithin}(\ell,-,-)$ and proved them reflexive and symmetric.  It
did not prove any statement relating two of them, and closed with the
observation that a Cauchy notion on this scale needs one.  This file supplies
it, in the form the completion will use:
\[
  \name{ScaleWithin}(\ctor{succ}\,\ell, a, b)
  \;\wedge\;
  \name{ScaleWithin}(\ctor{succ}\,\ell, b, c)
  \;\longrightarrow\;
  \name{ScaleWithin}(\ell, a, c).
\]
Two steps, each at one level finer than the target, compose into one step at
the target level.  This is the multiplicative triangle inequality, and the
level bookkeeping in it is not an artefact: by
\texttt{PrimeTensor/Scale.lean} the tolerance at level $\ctor{succ}\,\ell$ is
the square root of the tolerance at level $\ell$, so ``one level finer'' is
exactly ``half the error'' in the multiplicative sense, and two halves make a
whole.

Getting there needs three things the previous files left out: that ratios
compose, $\name{ratio}(a,c) = \name{ratio}(a,b)\cdot\name{ratio}(b,c)$; that
the squaring ladder is multiplicative,
$\name{scalePow}(xy,\ell) = \name{scalePow}(x,\ell)\,\name{scalePow}(y,\ell)$;
and the arithmetic core, that $x^2 < 2$ and $y^2 < 2$ force $xy < 2$.  The
third is where the work is, and it is done on the natural-number
cross-products underlying $\name{PrimeRatio}$ --- the file introduces no real
numbers and no square roots.

\subsection{Ratios compose}

\begin{lemma}[\lean{MulRat.inv\_mul}]\label{Scale:Laws:lem:inv-mul}
$a^{-1} \cdot a = 1$ for every $a : \name{MulRat}$.
\end{lemma}

\begin{leancode}
@[simp] theorem inv_mul (a : MulRat) : a⁻¹ * a = 1 := by
  rw [mul_comm]
  exact mul_inv a
\end{leancode}

\begin{proof}
Commute the product and apply $\name{mul\_inv}$ of
\texttt{PrimeTensor/MulRat.lean}.
\end{proof}

\begin{remark}[Filling a gap, not proving something new]
\label{Scale:Laws:rem:inv-mul}
\texttt{PrimeTensor/MulRat.lean} proved $a \cdot a^{-1} = 1$ and, separately,
commutativity, but never stated the left-handed form; since $\name{MulRat}$
carries a bare $\ctor{Mul}$ instance and not a group structure, nothing
derives it automatically.  It is marked $\ctor{simp}$ here because the
cancellation in Lemma~\ref{Scale:Laws:lem:ratio-comp} is where it is wanted.
\end{remark}

\begin{lemma}[\lean{MulRat.ratio\_comp}]\label{Scale:Laws:lem:ratio-comp}
$\name{ratio}(a,c) = \name{ratio}(a,b) \cdot \name{ratio}(b,c)$ for all
$a,b,c : \name{MulRat}$.
\end{lemma}

\begin{leancode}
theorem ratio_comp (a b c : MulRat) :
    ratio a c = ratio a b * ratio b c := by
  unfold ratio
  symm
  calc
    (a * b⁻¹) * (b * c⁻¹)
        = ((a * b⁻¹) * b) * c⁻¹ := (mul_assoc (a * b⁻¹) b c⁻¹).symm
    _ = (a * (b⁻¹ * b)) * c⁻¹ := by
      rw [mul_assoc a b⁻¹ b]
    _ = (a * 1) * c⁻¹ := by
      rw [inv_mul b]
    _ = a * c⁻¹ := by
      rw [mul_one]
\end{leancode}

\begin{proof}
Unfolding $\name{ratio}$ of \texttt{PrimeTensor/Order.lean}, the claim is
$a c^{-1} = (a b^{-1})(b c^{-1})$, and it is proved from right to left:
reassociate to expose $b^{-1} b$, cancel it by
Lemma~\ref{Scale:Laws:lem:inv-mul}, and remove the unit.  The intermediate
magnitude $b$ leaves no trace, which is the point --- the oriented ratio is a
cocycle.
\end{proof}

\subsection{The ladder is multiplicative}

\begin{lemma}[\lean{MulRat.mul\_mul\_shuffle}, private]
\label{Scale:Laws:lem:shuffle}
$(xy)(xy) = (xx)(yy)$ for all $x,y : \name{MulRat}$.
\end{lemma}

\begin{proof}
Six rewrites, alternating associativity and the single commutation
$yx = xy$:
\[
  (xy)(xy) = x(y(xy)) = x((yx)y) = x((xy)y) = (x(xy))y = ((xx)y)y = (xx)(yy).
\]
\end{proof}

\begin{remark}[Why this is spelled out]\label{Scale:Laws:rem:no-ac}
The identity is a pure consequence of associativity and commutativity, and on
a type carrying a $\ctor{CommMonoid}$ instance it would be one call to
$\ctor{ac\_rfl}$ or $\ctor{ring}$.  $\name{MulRat}$ carries only $\ctor{One}$,
$\ctor{Mul}$ and $\ctor{Inv}$ instances, with associativity and commutativity
available as named theorems and not as instance data, so the rewriting has to
be written out.  The contrast is visible inside this same file:
Theorem~\ref{Scale:Laws:thm:mul-lt-two} does its regrouping on $\mathbb{N}$,
where the algebraic instances are in place, and there each regrouping step is
a single $\ctor{ac\_rfl}$.
\end{remark}

\begin{lemma}[\lean{MulRat.scalePow\_mul}]\label{Scale:Laws:lem:scalepow-mul}
$\name{scalePow}(x \cdot y, \ell) =
 \name{scalePow}(x, \ell) \cdot \name{scalePow}(y, \ell)$
for all $x, y : \name{MulRat}$ and every level $\ell$.
\end{lemma}

\begin{leancode}
theorem scalePow_mul (x y : MulRat) :
    ∀ level : Depth,
      scalePow (x * y) level = scalePow x level * scalePow y level
  | .one => rfl
  | .succ level => by
      change
        scalePow (x * y) level * scalePow (x * y) level =
          (scalePow x level * scalePow x level) *
            (scalePow y level * scalePow y level)
      rw [scalePow_mul x y level]
      exact mul_mul_shuffle (scalePow x level) (scalePow y level)
\end{leancode}

\begin{proof}
By structural recursion on the level.  At $\ctor{one}$ both sides are $xy$, so
the claim is $\ctor{rfl}$.  At $\ctor{succ}\,\ell$ both sides unfold to
squares; the recursive call rewrites $\name{scalePow}(xy,\ell)$ to
$\name{scalePow}(x,\ell)\,\name{scalePow}(y,\ell)$, and
Lemma~\ref{Scale:Laws:lem:shuffle} regroups the square of that product into
the product of the squares.
\end{proof}

Read through the exponent formula of \texttt{PrimeTensor/Scale.lean}, this is
$(xy)^{n} = x^{n} y^{n}$ for $n = 2^{\size{\ell}-1}$; but it is proved
directly on the ladder, so no exponent arithmetic enters the file.

\subsection{The ruler as a natural number}

\begin{lemma}[\lean{MulRat.multiset\_eval\_pos}, private]
\label{Scale:Laws:lem:eval-pos}
$0 < q.\name{eval}$ for every $q : \name{PrimeMultiset}$.
\end{lemma}

\begin{proof}
By induction on the quotient, it suffices to treat a bag, and there
$\name{PrimeBag.one\_le\_eval}$ of \texttt{PrimeTensor/Prime.lean} gives
$1 \leq \name{eval}$.
\end{proof}

\begin{remark}[The same lemma, a third time]\label{Scale:Laws:rem:third-time}
This statement is proved privately in \texttt{PrimeTensor/MulRat.lean}, again
in \texttt{PrimeTensor/Order.lean} as $\name{multiset\_eval\_pos\_order}$, and
again here.  Each copy is $\ctor{private}$, so no later file can reach the
earlier one; the duplication is the price of that.  What it says is the
positivity that makes the cross-product order well behaved at all: every
barcode denotes a magnitude $\geq 1$, so no denominator can vanish and no
comparison degenerates.
\end{remark}

\begin{lemma}[\lean{MulRat.twoRatio\_upper\_eval} and
  \lean{twoRatio\_lower\_eval}]\label{Scale:Laws:lem:two-eval}
$\name{twoRatio}.\name{upper}.\name{eval} = 2$ and
$\name{twoRatio}.\name{lower}.\name{eval} = 1$.
\end{lemma}

\begin{leancode}
@[simp] theorem twoRatio_upper_eval : twoRatio.upper.eval = 2 := by
  norm_num [twoRatio, primeTwo, PrimeMultiset.eval_ofBag,
    PrimeBag.eval_factor, PrimeBag.eval_one]

@[simp] theorem twoRatio_lower_eval : twoRatio.lower.eval = 1 := by
  rfl
\end{leancode}

\begin{proof}
The numerator is the one-atom barcode carrying $\name{primeTwo}$: evaluating
through $\name{eval\_ofBag}$ and $\name{eval\_factor}$ gives
$\name{primeTwo}.\name{value} \cdot 1$, which $\ctor{norm\_num}$ reduces to
$2$.  The denominator is the pivot, whose evaluation is $1$ definitionally.
\end{proof}

These two are the hinge of the file.  Marked $\ctor{simp}$, they let every
comparison against the ruler be normalised into a statement about the literal
naturals $2$ and $1$, after which the remaining work is ordinary arithmetic on
$\mathbb{N}$ and nothing about barcodes is needed again.

\subsection{The arithmetic core}

\begin{theorem}[\lean{MulRat.mul\_lt\_two\_of\_sq\_lt\_two}]
\label{Scale:Laws:thm:mul-lt-two}
If $x \cdot x < \name{two}$ and $y \cdot y < \name{two}$ in $\name{MulRat}$,
then $x \cdot y < \name{two}$.
\end{theorem}

\begin{proof}
Both hypotheses are about classes, and $<$ on $\name{MulRat}$ was defined in
\texttt{PrimeTensor/Order.lean} by lifting $\name{PrimeRatio.Lt}$ along the
quotient, so choose representatives $x_r$ and $y_r$ and write
\[
  p = x_r.\name{upper}.\name{eval}, \quad
  q = x_r.\name{lower}.\name{eval}, \quad
  u = y_r.\name{upper}.\name{eval}, \quad
  v = y_r.\name{lower}.\name{eval},
\]
so that $x$ denotes $p/q$ and $y$ denotes $u/v$.  Products of barcodes
multiply componentwise, by $\name{upper\_mul}$ and $\name{lower\_mul}$ of
\texttt{PrimeTensor/Ratio.lean}, and
$\name{PrimeRatio.Lt}(a,b)$ is the cross-product inequality
$a.\name{upper}.\name{eval} \cdot b.\name{lower}.\name{eval} <
 b.\name{upper}.\name{eval} \cdot a.\name{lower}.\name{eval}$.  Normalising
the ruler with Lemma~\ref{Scale:Laws:lem:two-eval} turns the two hypotheses
and the goal into statements about naturals:
\[
  p^2 < 2q^2, \qquad u^2 < 2v^2, \qquad\text{goal:}\quad pu < 2qv .
\]
Squaring is strictly monotone on $\mathbb{N}$, so by
$\name{Nat.mul\_self\_lt\_mul\_self\_iff}$ it is enough to prove the squared
goal $(pu)^2 < (2qv)^2$, and that is a chain of two strict steps:
\[
  (pu)^2 = p^2 u^2
    \;<\; (2q^2)\,u^2
    \;=\; u^2\,(2q^2)
    \;<\; (2v^2)(2q^2)
    \;=\; (2qv)^2 .
\]
The first strict step multiplies the first hypothesis by $u^2$ and the second
multiplies the second hypothesis by $2q^2$; both need the multiplier to be
strictly positive, which is Lemma~\ref{Scale:Laws:lem:eval-pos} applied to
$y_r.\name{upper}$ and to $x_r.\name{lower}$.  The three equalities are
regroupings of a commutative product, each closed by $\ctor{ac\_rfl}$.
\end{proof}

\begin{designnote}[Squaring in place of a square root]
\label{Scale:Laws:note:no-sqrt}
Stated in the reals the theorem is immediate: $x^2 < 2$ and $y^2 < 2$ give
$x, y < \sqrt{2}$, hence $xy < 2$.  That argument is unavailable here for the
reason \texttt{PrimeTensor/Scale.lean} recorded --- $\sqrt{2}$ is not an
element of $\name{MulRat}$, which contains exactly the positive rationals, so
the intermediate bound cannot even be written down.

The proof above avoids it by never dividing the problem in half.  It multiplies
the two hypotheses together instead, obtaining $(xy)^2 < 4$, and then descends
from $(xy)^2 < 4 = 2^2$ to $xy < 2$ by strict monotonicity of squaring on
$\mathbb{N}$ --- a square root taken on the integers, where it is exact,
rather than on the magnitudes, where it does not exist.  Everything the file
needs from irrational tolerances is obtained this way: the tolerances are
never named, only their squares, and the squares are rational.
\end{designnote}

\subsection{The composition law}

\begin{theorem}[\lean{MulRat.scaleWithin\_comp}]
\label{Scale:Laws:thm:comp}
If $\name{ScaleWithin}(\ctor{succ}\,\ell, a, b)$ and
$\name{ScaleWithin}(\ctor{succ}\,\ell, b, c)$, then
$\name{ScaleWithin}(\ell, a, c)$.
\end{theorem}

\begin{leancode}
theorem scaleWithin_comp {level : Depth} {a b c : MulRat}
    (hab : ScaleWithin (.succ level) a b)
    (hbc : ScaleWithin (.succ level) b c) :
    ScaleWithin level a c := by
\end{leancode}

\begin{proof}
By definition $\name{scalePow}(r, \ctor{succ}\,\ell)$ is
$\name{scalePow}(r,\ell)^2$, so the four conjuncts of the two hypotheses read
\[
  \name{scalePow}(\name{ratio}(a,b),\ell)^2 < \name{two},
  \quad
  \name{scalePow}(\name{ratio}(b,a),\ell)^2 < \name{two},
\]
and the same pair for $(b,c)$ and $(c,b)$.

For the forward conjunct of the goal, rewrite
$\name{ratio}(a,c) = \name{ratio}(a,b)\,\name{ratio}(b,c)$ by
Lemma~\ref{Scale:Laws:lem:ratio-comp} and distribute the ladder over the
product by Lemma~\ref{Scale:Laws:lem:scalepow-mul}, reducing the goal to
\[
  \name{scalePow}(\name{ratio}(a,b),\ell) \cdot
  \name{scalePow}(\name{ratio}(b,c),\ell) < \name{two},
\]
which is Theorem~\ref{Scale:Laws:thm:mul-lt-two} applied to the two forward
hypotheses.  For the backward conjunct, rewrite instead
$\name{ratio}(c,a) = \name{ratio}(c,b)\,\name{ratio}(b,a)$ --- note the
reversed order --- and apply the theorem to the backward hypothesis for
$(c,b)$ and the backward hypothesis for $(b,a)$.
\end{proof}

\begin{remark}[The level cost is exactly right]\label{Scale:Laws:rem:cost}
By \texttt{PrimeTensor/Scale.lean} the tolerance at level $\ell$ is
$2^{1/2^{\size{\ell}-1}}$, and at $\ctor{succ}\,\ell$ it is
$2^{1/2^{\size{\ell}}}$, the square root of the former.  So the theorem says:
two ratios each below $t$ have product below $t^2$, with $t$ the finer
tolerance and $t^2$ the coarser one.  The inequality is sharp in the sense
that no smaller target level would follow --- $t \cdot t$ is exactly $t^2$,
not less --- so composing costs one level and cannot cost fewer.

That single-level cost is what makes the ladder usable downstream.  A Cauchy
sequence for this scale will be indexed by depth, and chaining $n$ steps of a
telescoping argument will consume $n$ levels; since the levels are the depths
of \texttt{PrimeTensor/Depth.lean} and those are unbounded, there is always
room to pay.
\end{remark}

\begin{remark}[Monotonicity in the level now follows]
\label{Scale:Laws:rem:mono}
\texttt{PrimeTensor/Scale.lean} noted that closeness at a finer level ought to
imply closeness at a coarser one, and that the argument needed a comparison
fact not then available.  Theorem~\ref{Scale:Laws:thm:mul-lt-two} is that
fact, and the implication is now a one-line consequence of
Theorem~\ref{Scale:Laws:thm:comp}: take the middle magnitude to be $a$ itself
and feed $\name{scaleWithin\_refl}(\ctor{succ}\,\ell, a)$ as the first
hypothesis, so that
\[
  \name{ScaleWithin}(\ctor{succ}\,\ell, a, b)
  \;\longrightarrow\;
  \name{ScaleWithin}(\ell, a, b).
\]
The file does not state this, and nothing below needs it in this form; it is
recorded here because it settles the question the previous file left open.
\end{remark}

\begin{remark}[What is still missing]\label{Scale:Laws:rem:missing}
The composition law is proved only in the shape it will be used ---
$\ctor{succ}\,\ell$ and $\ctor{succ}\,\ell$ into $\ell$.  There is no general
statement composing level $m$ with level $n$ into some level determined by
both, and no $n$-fold version chaining a whole finite path.  Nor is
$\name{ScaleWithin}$ shown to be an equivalence relation at any fixed level:
it is reflexive and symmetric, but transitivity at a single level is false in
general, and the level-shifted law above is what replaces it.  A completion
built on this scale has to be organised around that shift rather than around
an equivalence.
\end{remark}

\subsection*{Summary of the correspondence}

\begin{center}
\small
\begin{tabular}{@{}lll@{}}
\hline
Lean declaration & Kind & Here \\
\hline
\lean{MulRat.inv\_mul}                   & simp thm        & Lem.~\ref{Scale:Laws:lem:inv-mul} \\
\lean{MulRat.ratio\_comp}                & theorem         & Lem.~\ref{Scale:Laws:lem:ratio-comp} \\
\lean{MulRat.mul\_mul\_shuffle}          & private thm     & Lem.~\ref{Scale:Laws:lem:shuffle} \\
\lean{MulRat.scalePow\_mul}              & theorem         & Lem.~\ref{Scale:Laws:lem:scalepow-mul} \\
\lean{MulRat.multiset\_eval\_pos}        & private thm     & Lem.~\ref{Scale:Laws:lem:eval-pos} \\
\lean{MulRat.twoRatio\_upper\_eval}      & simp thm        & Lem.~\ref{Scale:Laws:lem:two-eval} \\
\lean{MulRat.twoRatio\_lower\_eval}      & simp thm        & Lem.~\ref{Scale:Laws:lem:two-eval} \\
\lean{MulRat.mul\_lt\_two\_of\_sq\_lt\_two} & theorem      & Thm.~\ref{Scale:Laws:thm:mul-lt-two} \\
\lean{MulRat.scaleWithin\_comp}          & theorem         & Thm.~\ref{Scale:Laws:thm:comp} \\
\hline
\end{tabular}
\end{center}

\noindent
Every declaration of \texttt{PrimeTensor/Scale/Laws.lean} appears above.  The
file contains no \lean{sorry}, no axiom, and no unproved named proposition.
Remarks~\ref{Scale:Laws:rem:cost}, \ref{Scale:Laws:rem:mono}
and~\ref{Scale:Laws:rem:missing} are stated for the reader and are not
declarations of the file.
